% `tex.web` §638 and §640: the `[N]` a TeX job prints as it works.
%
%     if term_offset>max_print_line-9 then print_ln
%     else if (term_offset>0)or(file_offset>0) then print_char(" ");
%     print_char("["); j:=9;
%     while (j>0)and(count(j)=0) do decr(j);
%     for k:=0 to j do
%       begin print_int(count(k)); if k<j then print_char("."); end;
%
% and §640 closes it — `if tracing_output<=0 then print_char("]")` — *after*
% the page has actually gone out, so anything raised while shipping falls
% between the two halves.
%
% It is `\count0` through the last nonzero one, dot-separated: `\count0` is
% the page number, and a document that uses `\count1` for a part number gets
% `[7.3]`. And it is written against the current column, so it joins whatever
% line is in progress with a space in front — which is why this waited on a
% printer that has a column at all. On a list of whole lines it would have
% gone on one of its own: the right information in the wrong place.
%
% With `\tracingoutput` positive the same numbers appear inside §638's
% `Completed box being shipped out [0]` announcement instead, which is a
% different line and was already right — tests/showbox/tracingoutput.tex.
\catcode`\{=1 \catcode`\}=2
\hsize=100pt \parindent=0pt \vsize=20pt
\font\r=cmr10 \r
\hbadness=10000 \vbadness=10000
\showthe\escapechar

% Two pages at the default counters.
A\par B\par

% ...and two more with `\count1` set, so the number has two parts.
\count0=7 \count1=3
C\par

% A `\message` leaves the line open, and the next `[N]` joins it.
\message{after}
D\par
\showthe\escapechar
\end
